\bitsection{背景与问题}{研究动机、必要概念与相关工作}

\begin{frame}{研究问题与汇报目标}
  \begin{block}{研究对象与应用条件}
    \bitplaceholder{明确研究对象、输入输出，以及需要满足的场景或物理约束。}
  \end{block}
  \begin{block}{现有方法的瓶颈}
    \begin{itemize}
      \item \bitplaceholder{说明具体困难及其产生原因。}
      \item \bitplaceholder{说明这个困难对性能或使用效果的影响。}
    \end{itemize}
  \end{block}
  \bitcallout{本次汇报：\bitplaceholder{讲清一个核心问题，并说明方法与证据。}}
  \bitsource{文献来源 / 项目背景：作者、年份、章节或链接。}
\end{frame}

\begin{frame}{基础概念与符号定义}
  \bitcolumns{%
    \begin{block}{概念定义}
      \bitplaceholder{先用一句话解释概念，再给出必要的数学定义。}
    \end{block}
    \[
      y_t=f_{\theta}(x_{\leq t},c)
    \]
    {\small\bitplaceholder{此公式仅作为排版占位，使用时替换。}}
  }{%
    \small\renewcommand{\arraystretch}{1.7}
    \begin{tabularx}{\linewidth}{@{}lX@{}}
      \toprule
      \bitstrong{符号} & \bitstrong{含义 / 单位}\\\midrule
      $x_{\leq t}$ & 历史输入或观测序列\\
      $c$ & 任务条件或外部信息\\
      $y_t$ & 当前预测或控制输出\\
      $\theta$ & 可学习参数\\\bottomrule
    \end{tabularx}
  }
  \bitcallout{直观理解：\bitplaceholder{用观众熟悉的具体场景解释符号。}}
\end{frame}

\begin{frame}{概念对照与问题转换}
  \bitcolumns{%
    \begin{block}{原始问题 / 方法 A}
      \bitimageplaceholder[.28\textheight]{原始问题示意图}
      \small\bitplaceholder{说明输入、输出和当前假设。}
    \end{block}
  }{%
    \begin{block}{新表述 / 方法 B}
      \bitimageplaceholder[.28\textheight]{转换后的问题示意图}
      \small\bitplaceholder{说明对应关系和新增能力。}
    \end{block}
  }
  \bitcallout{转换依据：\bitplaceholder{解释哪些量相互对应，为什么这样转换合理。}}
\end{frame}

\begin{frame}{相关工作与本文定位}
  \small\renewcommand{\arraystretch}{1.7}
  \begin{tabularx}{\textwidth}{@{}p{26mm}XX@{}}
    \toprule
    \bitstrong{方法类别} & \bitstrong{解决的问题} & \bitstrong{与本工作的关系}\\\midrule
    方法类别 A & \bitplaceholder{能力与适用条件} & \bitplaceholder{仍然存在的不足}\\
    方法类别 B & \bitplaceholder{能力与适用条件} & \bitplaceholder{可借鉴的设计}\\
    本文 / 本工作 & \bitplaceholder{主要改进} & \bitplaceholder{适用范围与前提}\\\bottomrule
  \end{tabularx}
  \bitcallout{工作定位：\bitplaceholder{相对最相关的方法，具体改变了什么。}}
  \bitsource{每类列出 1--2 篇代表作。论文汇报可在本页补充论文作者与发表信息。}
\end{frame}
